\chapter{1975 Nobel Prize in Economics}
\textbf{Laureates: }Leonid Kantorovich (USSR), Tjalling Koopmans (Netherlands/USA)

\section*{Reason for Award}
For their contributions to the theory of optimal allocation of resources

\section{What They Did}
Leonid Kantorovich pioneered linear programming, developing mathematical methods for determining the optimal allocation of limited resources to maximize output in industrial planning contexts. He created dual valuation methods that provided shadow prices for resources, enabling efficient planning in command economies. Tjalling Koopmans applied linear programming to transportation problems, minimizing shipping costs and resource allocation efficiency, establishing foundations of operations research and efficient resource distribution theory. Koopmans developed the concept of economic efficiency now bearing his name, defining Pareto optimality in rigorous mathematical terms. Together, their work established rigorous mathematical methods for solving resource allocation problems, enabling efficient planning and operations research to become essential tools in management science, logistics, manufacturing, and public policy.

\section{Impact on Economic Thinking}
Together, their work established rigorous mathematical methods for solving resource allocation problems, enabling efficient planning in command economies and operations research became an essential tool for management science, logistics, manufacturing, and public policy. Optimization techniques provided a systematic approach to scheduling production, routing vehicles, allocating budgets, and designing networks, achieving maximum efficiency at minimum cost. Linear programming, integer programming, and branch-and-bound algorithms derive from their foundational contributions. Their work demonstrated that mathematical optimization could solve practical economic problems, bridging the gap between theoretical economics and real-world decision-making. The duality theory they developed provided economic interpretation of mathematical constraints, revealing shadow prices as marginal values of resources.

\section{Modern Perspectives Today}
Linear programming and operations research remain fundamental to computer science, engineering, logistics, supply chain management, and optimization theory. The field has extended into integer programming, nonlinear programming, and combinatorial optimization, finding wide application in telecommunications, healthcare, finance, and machine learning. Operations research methods are increasingly important for managing complex systems and optimizing resource allocation under uncertainty. Algorithmic game theory extends these principles to strategic settings with competing agents. Modern applications include network design, scheduling, inventory management, and portfolio optimization. The mathematical foundations laid by Kantorovich and Koopmans continue to underpin optimization theory and its applications across economics, engineering, and computer science. Their work remains essential for understanding efficient resource allocation in both planned and market economies.
